\documentclass[12pt, a4paper]{article}
\usepackage[utf8]{inputenc}
\usepackage[spanish]{babel}
\usepackage{amsmath}
\usepackage{amssymb}
\usepackage{booktabs}
\usepackage{geometry}
\usepackage{xcolor}
\usepackage{graphicx}

\geometry{margin=2.5cm}

\title{\textbf{Prueba Calc}\\[0.5em]
\large LaTeX + celdas Python con \texttt{\textbackslash py\{\}}}
\author{}
\date{\today}

\begin{document}
\maketitle
\tableofcontents
\newpage

% ===========================================================
\section{Cálculos básicos}
% ===========================================================

%#python
import math

# Variables básicas
radio       = 5.0
altura      = 10.0
gravedad    = 9.81   # m/s²
masa        = 70.0   # kg

# Cálculos
area_circulo    = math.pi * radio**2
perimetro       = 2 * math.pi * radio
vol_cilindro    = area_circulo * altura
peso            = masa * gravedad

print(f"Área círculo  : {area_circulo:.4f} m²")
print(f"Volumen cil.  : {vol_cilindro:.4f} m³")
print(f"Peso          : {peso:.2f} N")
%#end

Para un círculo de radio $r = \py{radio}$ m:

\[
  A = \pi r^2 = \py{f"{area_circulo:.4f}"} \text{ m}^2
  \qquad
  P = 2\pi r = \py{f"{perimetro:.4f}"} \text{ m}
\]

El volumen del cilindro de altura $h = \py{altura}$ m es:

\[
  V = A \cdot h = \py{f"{area_circulo:.4f}"} \times \py{altura}
    = \py{f"{vol_cilindro:.4f}"} \text{ m}^3
\]

Una persona de \py{masa} kg pesa $\py{f"{peso:.1f}"}$ N
($\py{f"{peso/1000:.3f}"}$ kN). También puedes operar directamente
en el documento: $2A - \tfrac{3}{2} = \py{area_circulo*2 - 3/2}$.


% ===========================================================
\section{Secuencias y tablas}
% ===========================================================

%#python
# Tabla de potencias (devuelve filas LaTeX como cadena)
filas_potencias = "\n    ".join(
    f"{n} & {n**2} & {n**3} & {n**4} \\\\"
    for n in range(1, 8)
)

# Números de Fibonacci
def fib(n):
    a, b = 0, 1
    seq = []
    for _ in range(n):
        seq.append(a)
        a, b = b, a + b
    return seq

fib10 = fib(10)
fib_str = ", ".join(str(x) for x in fib10)
fib_sum = sum(fib10)

print("Fibonacci (10):", fib_str)
print("Suma:", fib_sum)
%#end

\subsection{Potencias}

\begin{table}[h]
  \centering
  \caption{Potencias de $n$ calculadas en Python}
  \begin{tabular}{c c c c}
    \toprule
    $n$ & $n^2$ & $n^3$ & $n^4$ \\
    \midrule
    \py{filas_potencias}
    \bottomrule
  \end{tabular}
\end{table}

\subsection{Sucesión de Fibonacci}

Los primeros 10 términos de la sucesión de Fibonacci son:

\[
  \py{fib_str}
\]

Su suma es $\py{fib_sum}$.


% ===========================================================
\section{Trigonometría}
% ===========================================================

%#python
import math

angulos_deg = [0, 30, 45, 60, 90]

# Tabla trigonométrica (notación matemática: ^{\circ} para grados)
filas_trig = ""
for deg in angulos_deg:
    rad = math.radians(deg)
    s   = round(math.sin(rad), 4)
    c   = round(math.cos(rad), 4)
    t   = r"\infty" if deg == 90 else str(round(math.tan(rad), 4))
    filas_trig += f"    ${deg}^{{\\circ}}$ & ${s}$ & ${c}$ & ${t}$ \\\\\n"

# Identidad de Euler evaluada en pi
euler_re = round(math.cos(math.pi), 6)   # = -1
euler_im = round(math.sin(math.pi), 6)   # ≈ 0

print("Tabla trig. generada.")
print(f"e^(iπ) = {euler_re} + {euler_im}i")
%#end

\subsection{Tabla trigonométrica}

\begin{table}[h]
  \centering
  \caption{Valores de seno, coseno y tangente}
  \begin{tabular}{c c c c}
    \toprule
    Ángulo & $\sin$ & $\cos$ & $\tan$ \\
    \midrule
\py{filas_trig}
    \bottomrule
  \end{tabular}
\end{table}

\subsection{Identidad de Euler}

La famosa identidad $e^{i\pi} + 1 = 0$ verificada numéricamente:

\[
  e^{i\pi} = \cos(\pi) + i\sin(\pi)
           = \py{euler_re} + \py{euler_im}\,i
\]


% ===========================================================
\section{Estadística}
% ===========================================================

%#python
import math

datos = [12.3, 15.7, 9.8, 22.1, 18.4, 11.2, 25.6, 14.9, 17.3, 20.0]
n     = len(datos)
media = sum(datos) / n
varianza = sum((x - media)**2 for x in datos) / n
desv_std = math.sqrt(varianza)
minimo   = min(datos)
maximo   = max(datos)
rango    = maximo - minimo
datos_ord = sorted(datos)
mediana = (datos_ord[n//2 - 1] + datos_ord[n//2]) / 2

datos_str = ", ".join(f"{x}" for x in datos)

print(f"N={n}  media={media:.3f}  σ={desv_std:.3f}  mediana={mediana:.1f}")
%#end

Para el conjunto de datos $\{\ \py{datos_str}\ \}$ ($n = \py{n}$):

\[
  \bar{x} = \py{f"{media:.3f}"}
  \qquad
  \sigma  = \py{f"{desv_std:.3f}"}
  \qquad
  \tilde{x} = \py{f"{mediana:.1f}"}
\]

Rango: $[\py{minimo},\ \py{maximo}]$, amplitud $= \py{f"{rango:.1f}"}$.


% ===========================================================
\section{Expresiones directas en el texto}
% ===========================================================

También puedes escribir expresiones Python directamente en el texto
sin declarar nada en una celda:

\begin{itemize}
  \item $\sqrt{2} = \py{round(math.sqrt(2), 6)}$
  \item $\ln(e^3) = \py{round(math.log(math.e**3), 6)}$
  \item $100! \mod 97 = \py{math.factorial(100) % 97}$
  \item Longitud de la lista anterior: \py{len(datos)} elementos
\end{itemize}

\end{document}
